\subsection{Robustness}\label{sec:robust}
\emph{Warm-starting NSGA-II.} QHP-Net reuses its previous population; NSGA-II re-optimises from scratch. Giving NSGA-II the same warm start does not help it (FND 1,881 vs.\ 1,884 lossy, 1,968 vs.\ 1,965 ideal, both $p=0.48$), and QHP-Net still leads by $+252$ [$+241,+264$] rounds ($d_z=8.5$). The gain comes from the objective, not the memory.

\emph{Heterogeneous energy.} With three energy levels (20\% of nodes advanced at $2E_0$, 10\% super at $3E_0$), QHP-Net's FND rises to 2,894, and its lead over NSGA-II grows to $+321$ rounds (Table~\ref{tab:robust}). Because the fitness divides drain by residual energy, load moves onto high-energy nodes without any heterogeneity-specific rule. (DQN was trained on homogeneous fields only.)

\emph{Charging the uplink.} The headline runs, like the earlier study, let centralised protocols learn node state free from data traffic. If every alive node instead pays a 200-bit report at each re-clustering, QHP-Net loses 38 rounds and NSGA-II 43; Fuzzy, DQN and HIECF lose 98--233, because they re-cluster every round. QHP-Net still wins all 30 topologies, by at least 200 rounds.

\emph{Network size and field.} Table~\ref{tab:sweep} varies $N\in\{50,100,150\}$ and the field side $W$ from 50 to 150\,m (sink at $(W/2,1.5W)$, 15 paired topologies per cell). QHP-Net has the highest mean FND in all nine cells, beats the best other protocol in each under a paired $t$-test ($p<3\times10^{-6}$), and wins 134 of 135 topologies; the one loss, by 12 rounds, is 50 nodes on a 150\,m field. The margin grows with density: $+185$ to $+274$ rounds at 100\,m and $+299$ to $+809$ at 150\,m, where sink hops pay the multipath amplifier. On the 50\,m field, where Direct is competitive (2,297 at $N{=}50$), the direct regime uses it selectively and still adds 400--542 rounds.

\begin{table}[t]
\caption{Robustness on the lossy channel, 30 paired runs: three-level heterogeneity (top) and uplink state reports charged to every centralised protocol (bottom). AUC in $10^3$ node-rounds.}
\label{tab:robust}
\centering\footnotesize
\setlength{\tabcolsep}{2.6pt}
\begin{tabular}{@{}lrrrr@{}}
\toprule
Protocol & FND & HND & AUC & Yield\\
\midrule
\multicolumn{5}{@{}l}{\emph{Heterogeneous} ($m=0.2$, $b=0.1$)}\\
LEACH & 1,057 $\pm$ 47 & 1,797 $\pm$ 31 & 225.3 & 0.877 \\
PEGASIS & 1,673 $\pm$ 280 & 2,227 $\pm$ 5 & 295.1 & 0.834 \\
NSGA-II & 2,573 $\pm$ 81 & 2,688 $\pm$ 65 & 280.5 & 0.888 \\
Fuzzy-T2 & 2,085 $\pm$ 89 & 2,690 $\pm$ 59 & 265.8 & 0.898 \\
DQN & 1,506 $\pm$ 86 & 2,181 $\pm$ 44 & 258.6 & 0.878 \\
HIECF & 2,092 $\pm$ 162 & 2,654 $\pm$ 73 & 266.5 & 0.905 \\
QHP-Net & 2,894 $\pm$ 64 & 2,963 $\pm$ 22 & 296.1 & 0.983 \\

\midrule
\multicolumn{5}{@{}l}{\emph{Uplink charged}}\\
NSGA-II & 1,841 $\pm$ 32 & 1,942 $\pm$ 29 & 197.4 & 0.889 \\
Fuzzy-T2 & 1,366 $\pm$ 55 & 1,761 $\pm$ 33 & 173.3 & 0.915 \\
DQN & 1,170 $\pm$ 70 & 1,696 $\pm$ 42 & 172.9 & 0.883 \\
HIECF & 1,392 $\pm$ 112 & 1,772 $\pm$ 38 & 174.2 & 0.924 \\
QHP-Net & 2,095 $\pm$ 18 & 2,131 $\pm$ 17 & 213.1 & 0.983 \\

\bottomrule
\end{tabular}
\end{table}

\begin{table}[t]
\caption{Mean FND over 15 paired topologies per cell (lossy). Columns: field side $W$ (m) / number of nodes $N$.}
\label{tab:sweep}
\centering
\resizebox{\columnwidth}{!}{%
\begin{tabular}{@{}lrrrrrrrrr@{}}
\toprule
& \multicolumn{3}{c}{$W=50$} & \multicolumn{3}{c}{$W=100$} & \multicolumn{3}{c}{$W=150$}\\
\cmidrule(lr){2-4}\cmidrule(lr){5-7}\cmidrule(l){8-10}
Protocol & 50 & 100 & 150 & 50 & 100 & 150 & 50 & 100 & 150\\
\midrule
Direct & 2297 & 2020 & 1976 & 122 & 114 & 112 & 25 & 23 & 23 \\
LEACH & 1808 & 1620 & 1578 & 642 & 1045 & 1106 & 154 & 313 & 354 \\
PEGASIS & 2090 & 2083 & 2080 & 1463 & 1626 & 1554 & 624 & 772 & 810 \\
NSGA-II & 2130 & 2186 & 2176 & 1787 & 1883 & 1875 & 761 & 1107 & 1124 \\
DQN & 1562 & 1614 & 1677 & 1322 & 1389 & 1454 & 447 & 605 & 673 \\
HIECF & 1637 & 1651 & 1643 & 1422 & 1478 & 1491 & 412 & 500 & 569 \\
\textbf{QHP-Net} & 2697 & 2723 & 2718 & 1972 & 2138 & 2149 & 1060 & 1826 & 1933 \\

\bottomrule
\end{tabular}}
\end{table}
